\chapter*{Contrato metodológico global}
\addcontentsline{toc}{chapter}{Contrato metodológico global}
\label{chap:contract}
\statusbox{Regla de oro}{
Las 59 etiquetas ocultas de FIRA nunca se cargan, reconstruyen indirectamente ni se utilizan
para selección. Toda operación que usa $y$ se ajusta exclusivamente con etiquetas visibles
dentro de la partición correspondiente. Las operaciones estrictamente $X$-only pueden usar
las covariables de las 197 parcelas cuando el protocolo transductivo lo permite.
}

\section*{Universo estadístico}
Sea $\Train$ el conjunto de 138 parcelas con rendimiento observado y $\Target$ el conjunto
de 59 parcelas de predicción. Para cada parcela $i$ se observa $x_i\in\R^p$ y, si
$i\in\Train$, un rendimiento $y_i$. El objeto de inferencia es
\[
\widehat{\bm y}_{\Target}=(\hat y_i:i\in\Target),
\]
comparado por el evaluador con el vector reservado $\bm y_{\Target}$.

La información disponible durante el desarrollo es
\[
\mathcal I=
\{X_{\Train},\bm y_{\Train},X_{\Target},G,E_{\mathrm{public}}\},
\]
donde $G$ denota geometría/estructura espacial y $E_{\mathrm{public}}$ denota evidencia
externa pública permitida. Por construcción,
\[
\bm y_{\Target}\notin\mathcal I.
\]

\section*{Inducción, transducción y por qué la distinción importa}
En un problema inductivo convencional se pretende aprender una función
$f:\mathcal X\rightarrow\mathbb R$ que generalice a puntos futuros no conocidos. Aquí el
conjunto de consultas $X_{\Target}$ está fijado de antemano. Esto habilita transformaciones
transductivas dependientes de la geometría conjunta de $\Xall$: estandarización $X$-only,
PCA $X$-only, densidad, vecindarios, grafos y métricas de soporte. La restricción es
\[
T_X=T(X_{\Train},X_{\Target})\quad\text{es admisible si no usa }\bm y_{\Target},
\]
mientras cualquier transformación supervisada debe recalcularse usando exclusivamente las
etiquetas permitidas por el fold.

\section*{Métricas}
Para un conjunto de evaluación $S$,
\begin{align}
\RMSE(S)&=\sqrt{\frac{1}{|S|}\sum_{i\in S}(y_i-\hat y_i)^2},\\
\MAE(S)&=\frac{1}{|S|}\sum_{i\in S}|y_i-\hat y_i|,\\
R^2(S)&=
1-\frac{\sum_{i\in S}(y_i-\hat y_i)^2}
{\sum_{i\in S}(y_i-\bar y_S)^2}.
\end{align}

La selección principal del Checkpoint~04 usa RMSE medio entre pseudo-competiciones
target-matched. También se conservan RMSE pooled, MAE, peor split y protocolos de estrés,
porque con $n=138$ las conclusiones pueden cambiar si un split espacial domina el promedio.

\section*{Jerarquía de evidencia}
La interpretación final sigue esta prioridad:
\begin{enumerate}
\item invariantes y semántica de fuentes;
\item predicciones pseudo-target realmente out-of-sample respecto a $y$;
\item estabilidad entre splits y protocolos;
\item robustez espacial/grupada;
\item diagnósticos de soporte, discrepancia y plausibilidad;
\item sólo al final, predicciones de las 59 parcelas, sin utilizarlas para ajustar reglas.
\end{enumerate}
Esta jerarquía evita seleccionar modelos por el aspecto visual de sus 59 valores finales,
optimizar hiperparámetros contra un único protocolo conveniente o reinterpretar información
pública municipal como si fuese etiqueta parcelaria.
